\section{تحلیل اکتشافی داده‌ها}

پیش از آموزش طبقه‌بند، یک تحلیل اکتشافی داده‌ها (\lr{Exploratory Data Analysis, EDA}) انجام شد تا
درک بهتری از توزیع ویژگی‌ها، وجود مقادیر پرت و روابط بین متغیرها به دست آید. تمامی نمودارها و
تصویرها با استفاده از کتابخانه‌های \lr{matplotlib} و \lr{seaborn} در پایتون تولید شده‌اند.

\subsection{هیستوگرام ویژگی‌های عددی}

برای بررسی اولیهٔ ویژگی‌های عددی، هیستوگرام‌هایی برای پنج ویژگی پیوستهٔ اصلی یعنی
\lr{age}، \lr{chol}، \lr{oldpeak}، \lr{thalach} و \lr{trestbps} رسم شد. این نمودارها شکل کلی توزیع،
میزان کجی، و فراوانی مقادیر حدی را نشان می‌دهند. در ادامه، الگوهای مشاهده‌شده در هر ویژگی
به‌طور خلاصه توصیف می‌شوند.

\paragraph{سن (\lr{age}).}
نمودار شکل \ref{fig:hist-age} نشان می‌دهد که سن بیماران تقریباً در بازهٔ ۳۰ تا ۸۰ سال قرار دارد.
توزیع سن شکلی نزدیک به نرمال با قله‌ای مشخص در حدود ۵۵ تا ۶۰ سال دارد. بخش عمدهٔ نمونه‌ها را
بزرگسالان میانسال و سالمند تشکیل می‌دهند و تعداد افراد زیر ۴۰ سال بسیار کمتر است. این الگو از
نظر بالینی طبیعی است، زیرا ریسک بیماری قلبی با افزایش سن به‌طور قابل توجهی بیشتر می‌شود. پس از
سنین حدود ۶۵ تا ۷۰ سال، کاهش ملایمی در فراوانی نمونه‌ها دیده می‌شود که منعکس‌کنندهٔ تعداد کمتر
بیماران بسیار سالمند در این مجموعه‌داده است.

\begin{figure}[H]
	\centering
	\includegraphics[width=0.7\textwidth]{hist_age.png}
	\caption{هیستوگرام ویژگی \lr{age}.}
	\label{fig:hist-age}
\end{figure}

\paragraph{کلسترول (\lr{chol}).}
هیستوگرام \lr{chol} (شکل \ref{fig:hist-chol}) توزیعی با کجی ملایم به سمت راست را نشان می‌دهد.
بیشتر مقادیر در بازهٔ ۲۰۰ تا ۲۶۰ میلی‌گرم بر دسی‌لیتر قرار دارند؛ این محدوده در عمل بالینی با
سطوح مرزی یا بالا برای کلسترول سازگار است. دمی بلند در سمت راست تا حدود ۴۰۰ تا ۵۵۰ میلی‌گرم
بر دسی‌لیتر ادامه می‌یابد که نشان‌دهندهٔ تعداد اندکی از بیماران با اختلالات شدید چربی خون است.
این مقادیر پرت از نظر پزشکی قابل قبول و با عوامل خطر شناخته‌شدهٔ بیماری قلبی هم‌خوان هستند.

\begin{figure}[H]
	\centering
	\includegraphics[width=0.7\textwidth]{hist_chol.png}
	\caption{هیستوگرام ویژگی \lr{chol}.}
	\label{fig:hist-chol}
\end{figure}

\paragraph{\lr{oldpeak}.}
ویژگی \lr{oldpeak} که تغییر قطعهٔ \lr{ST} ناشی از ورزش را نشان می‌دهد، در میان همهٔ ویژگی‌های
عددی، بیشترین کجی به سمت راست را دارد (شکل \ref{fig:hist-oldpeak}). بخش بزرگی از بیماران دارای
مقداری نزدیک به صفر هستند که نشان‌دهندهٔ تغییر ناچیز در قطعهٔ \lr{ST} است. با افزایش مقدار
\lr{oldpeak}، فراوانی به‌سرعت کاهش می‌یابد و تنها تعداد اندکی نمونه با مقادیر بالاتر از ۳ و
مقادیر بسیار کمی با مقادیر حدود ۵ یا ۶ دیده می‌شود. این الگو نیز با انتظارات بالینی هم‌خوانی دارد،
زیرا افت شدید \lr{ST} در جمعیت کلی بیماران قلبی پدیده‌ای نسبتاً نادر است.

\begin{figure}[H]
	\centering
	\includegraphics[width=0.7\textwidth]{hist_oldpeak.png}
	\caption{هیستوگرام ویژگی \lr{oldpeak}.}
	\label{fig:hist-oldpeak}
\end{figure}

\paragraph{\lr{thalach}.}
حداکثر ضربان قلب (\lr{thalach}) توزیعی تقریباً متقارن و زنگوله‌ای‌شکل دارد
(شکل \ref{fig:hist-thalach}). بیشترین فراوانی در بازهٔ ۱۵۰ تا ۱۷۰ ضربان در دقیقه در طول تست
ورزشی مشاهده می‌شود. مقادیر بسیار پایین (کمتر از ۱۰۰) بسیار نادر هستند و می‌توانند نشان‌دهندهٔ
مشکلات شدید قلبی یا ناتمام‌ماندن تست ورزش باشند. این رفتار با مشاهدات بالینی روی جمعیت‌های
مورد ارزیابی در تست‌های ورزش هم‌خوانی خوبی دارد.

\begin{figure}[H]
	\centering
	\includegraphics[width=0.7\textwidth]{hist_thalach.png}
	\caption{هیستوگرام ویژگی \lr{thalach}.}
	\label{fig:hist-thalach}
\end{figure}

\paragraph{فشار خون استراحت (\lr{trestbps}).}
همان‌طور که در شکل \ref{fig:hist-trestbps} دیده می‌شود، توزیع \lr{trestbps} اندکی به سمت مقادیر
بالاتر متمایل است. بیشتر مشاهدات در بازهٔ ۱۲۰ تا ۱۴۰ میلی‌متر جیوه قرار دارند که به فشار خون
خفیف تا متوسط (پرفشاری خون) مربوط است و یک عامل خطر شایع برای بیماری قلبی محسوب می‌شود. مقادیر
بالاتر، تا حدود ۱۸۰ تا ۲۰۰ میلی‌متر جیوه، با فراوانی کم‌تر مشاهده می‌شوند، اما از نظر بالینی
برای بیماران پرخطر یا درمان‌نشده کاملاً واقع‌بینانه‌اند.

\begin{figure}[H]
	\centering
	\includegraphics[width=0.7\textwidth]{hist_trestbps.png}
	\caption{هیستوگرام ویژگی \lr{trestbps}.}
	\label{fig:hist-trestbps}
\end{figure}

در مجموع، هیستوگرام‌ها نشان می‌دهند که ویژگی‌های عددی ترکیبی از توزیع‌های تقریباً نرمال،
با کجی ملایم و همچنین توزیع‌های شدیداً کج را تشکیل می‌دهند که همگی با رفتار واقعی داده‌های
پزشکی سازگارند. این نتایج ضرورت نرمال‌سازی ویژگی‌ها را پیش از به‌کارگیری الگوریتم‌های مبتنی
بر فاصله‌ای مانند \lr{KNN} تقویت می‌کند.

\subsection{نمودارهای جعبه‌ای و مقادیر پرت}

برای بررسی پراکندگی، گرایش مرکزی و مقادیر پرت، نمودارهای جعبه‌ای برای پنج ویژگی عددی اصلی رسم شد.
این نمودارها دامنهٔ تغییرات هر ویژگی، موقعیت میانه و وجود مقادیر غیرعادی را به‌خوبی نمایش می‌دهند
و درک بهتری از رفتار اندازه‌گیری‌های بالینی فراهم می‌کنند.

\paragraph{سن (\lr{age}).}
نمودار جعبه‌ای \lr{age} (شکل \ref{fig:boxplot-age}) نشان می‌دهد که اغلب بیماران در بازهٔ تقریبی
۵۰ تا ۶۲ سال قرار دارند و میانه در حدود ۵۶ تا ۵۸ سال است. مقادیر پرت قابل‌توجهی مشاهده نمی‌شود
و توزیع سنی نسبتاً یکنواخت و متمرکزی دارد. این الگو با انتظار بالینی مبنی بر شیوع بالاتر بیماری
قلبی در بزرگسالان میانسال و سالمند سازگار است.

\begin{figure}[H]
	\centering
	\includegraphics[width=0.7\textwidth]{boxplot_age.png}
	\caption{نمودار جعبه‌ای ویژگی \lr{age}.}
	\label{fig:boxplot-age}
\end{figure}

\paragraph{کلسترول (\lr{chol}).}
نمودار جعبه‌ای \lr{chol} (شکل \ref{fig:boxplot-chol}) وجود تعداد قابل‌توجهی مقادیر پرت در سمت
بالاتر را نشان می‌دهد؛ مقادیری بالاتر از ۴۰۰ و حتی نزدیک ۵۵۰ میلی‌گرم بر دسی‌لیتر. دامنهٔ بین‌چارکی
تقریباً بین ۲۰۰ تا ۲۷۵ میلی‌گرم بر دسی‌لیتر قرار دارد و میانه در حدود ۲۴۰ تا ۲۵۰ است. این مقادیر
پرت از نظر بالینی منطقی بوده و می‌توانند بیانگر موارد شدید هیپرکلسترولمی باشند که از عوامل خطر
مهم برای بیماری‌های قلبی-عروقی هستند.

\begin{figure}[H]
	\centering
	\includegraphics[width=0.7\textwidth]{boxplot_chol.png}
	\caption{نمودار جعبه‌ای ویژگی \lr{chol}.}
	\label{fig:boxplot-chol}
\end{figure}

\paragraph{\lr{oldpeak}.}
نمودار جعبه‌ای \lr{oldpeak} (شکل \ref{fig:boxplot-oldpeak}) توزیعی با کجی شدید به سمت راست
را تأیید می‌کند؛ به‌طوری که بیشترین بخش مقادیر بین ۰ تا حدود ۱.۵ قرار دارد. چند مقدار پرت واضح
با مقادیر بالاتر از ۵ نیز دیده می‌شود که بیانگر موارد نادر ولی از نظر بالینی بسیار قابل‌توجه
افت شدید \lr{ST} در حین تست ورزش هستند. این مقادیر به ساختار دنبالهٔ سنگین در هیستوگرام این
ویژگی کمک می‌کنند.

\begin{figure}[H]
	\centering
	\includegraphics[width=0.7\textwidth]{boxplot_oldpeak.png}
	\caption{نمودار جعبه‌ای ویژگی \lr{oldpeak}.}
	\label{fig:boxplot-oldpeak}
\end{figure}

\paragraph{\lr{thalach}.}
در نمودار جعبه‌ای \lr{thalach} (شکل \ref{fig:boxplot-thalach}) بیشتر داده‌ها بین حدود
۱۴۰ تا ۱۷۰ ضربان در دقیقه قرار دارند و میانه در حدود ۱۵۵ است. یک مقدار پرت پایین در حدود ۸۰
ضربان در دقیقه مشاهده می‌شود که ممکن است به محدودیت شدید عملکرد قلب یا ناقص‌بودن تست ورزش مربوط
باشد. مقادیر بالاتر نیز تا حدود ۲۰۰ ضربان در دقیقه امتداد دارند که در زمینهٔ تست ورزش قابل‌قبول
و از نظر فیزیولوژیک محتمل هستند.

\begin{figure}[H]
	\centering
	\includegraphics[width=0.7\textwidth]{boxplot_thalach.png}
	\caption{نمودار جعبه‌ای ویژگی \lr{thalach}.}
	\label{fig:boxplot-thalach}
\end{figure}

\paragraph{فشار خون استراحت (\lr{trestbps}).}
نمودار جعبه‌ای \lr{trestbps} (شکل \ref{fig:boxplot-trestbps}) توزیعی معمول در محیط‌های
بالینی را نشان می‌دهد؛ به‌طوری که بیشتر مقادیر در بازهٔ ۱۲۰ تا ۱۴۰ میلی‌متر جیوه و میانه
در حدود ۱۳۰ تا ۱۳۵ قرار دارد. چند مقدار پرت چشمگیر در محدودهٔ ۱۷۰ تا ۲۰۰ میلی‌متر جیوه
وجود دارد که می‌تواند نمایانگر موارد پرفشاری خون شدید باشد. این الگو با الگوهای شناخته‌شدهٔ
ریسک در بیماران قلبی سازگار است.

\begin{figure}[H]
	\centering
	\includegraphics[width=0.7\textwidth]{boxplot_trestbps.png}
	\caption{نمودار جعبه‌ای ویژگی \lr{trestbps}.}
	\label{fig:boxplot-trestbps}
\end{figure}

به‌طور کلی، نمودارهای جعبه‌ای حضور مقادیر پرت معنادار در ویژگی‌هایی مانند \lr{chol}،
\lr{oldpeak} و \lr{trestbps} را تأیید می‌کنند، در حالی که ویژگی‌هایی مانند \lr{age} و
\lr{thalach} توزیع پایدارتر و متمرکزتری دارند. با توجه به این‌که الگوریتم \lr{KNN} نسبت به
تعداد کمی از مقادیر پرت جداافتاده (به‌ویژه پس از نرمال‌سازی) نسبتاً مقاوم است، این مقادیر
در این پروژه حذف نشدند.

\subsection{نمودارهای \lr{Q--Q}}

برای بررسی میزان نزدیکی ویژگی‌های عددی به توزیع نرمال، نمودارهای \lr{Q--Q} برای سه ویژگی
نماینده یعنی \lr{age}، \lr{chol} و \lr{thalach} رسم شد. در هر نمودار، مقادیر مرتب‌شدهٔ نمونه
با کوانتیل‌های یک توزیع نرمال با میانگین و واریانس مشابه مقایسه می‌شوند. انحراف از خط قطری،
نشان‌دهندهٔ فاصله از نرمال بودن است.

\paragraph{سن (\lr{age}).}
نمودار \lr{Q--Q} مربوط به \lr{age} (شکل \ref{fig:qq-age}) نشان می‌دهد که بخش عمدهٔ نقاط
در نزدیکی خط قطری قرار دارند؛ این موضوع نشان‌دهندهٔ نزدیک بودن بخش مرکزی توزیع به نرمال است.
با این حال، در دو دم توزیع انحراف‌های قابل‌توجهی دیده می‌شود: سنین پایین‌تر اندکی زیر خط
مرجع و سنین بالاتر (حدوداً بالای ۷۵ سال) کمی بالای آن قرار می‌گیرند. این الگو بیانگر دنباله‌های
نسبتاً سنگین اما ساختار کلی نسبتاً متقارن است؛ رفتاری که برای متغیر سن در داده‌های بالینی
تا حد زیادی معمول است.

\begin{figure}[H]
	\centering
	\includegraphics[width=0.7\textwidth]{qq_age.png}
	\caption{نمودار \lr{Q--Q} برای ویژگی \lr{age}.}
	\label{fig:qq-age}
\end{figure}

\paragraph{کلسترول (\lr{chol}).}
نمودار \lr{Q--Q} برای \lr{chol} (شکل \ref{fig:qq-chol}) انحراف‌های بسیار شدیدتری را نشان می‌دهد،
به‌ویژه در دم بالایی توزیع. در حالی که بخش مرکزی داده‌ها تا حدی با خط قطری هم‌راستا است، چندین
مقدار بسیار بالا (بالاتر از ۴۰۰ میلی‌گرم بر دسی‌لیتر) به‌وضوح بالاتر از کوانتیل‌های نرمال متناظر
قرار می‌گیرند. این نقاط متناظر با موارد شدید هیپرکلسترولمی هستند و رفتار کج‌به‌راست توزیع \lr{chol}
را تأیید می‌کنند. در نتیجه، این ویژگی با توزیع نرمال فاصلهٔ قابل‌توجهی دارد و حاوی مقادیر پرت
بالینی مهم است.

\begin{figure}[H]
	\centering
	\includegraphics[width=0.7\textwidth]{qq_chol.png}
	\caption{نمودار \lr{Q--Q} برای ویژگی \lr{chol}.}
	\label{fig:qq-chol}
\end{figure}

\paragraph{\lr{thalach}.}
نمودار \lr{Q--Q} ویژگی \lr{thalach} (شکل \ref{fig:qq-thalach}) نشان می‌دهد که بخش مرکزی توزیع
تقریباً روی خط قطری قرار دارد و رفتاری نزدیک به نرمال را نمایش می‌دهد. با این حال، در دم پایین،
چند مقدار بسیار کم (در حدود ۷۰ تا ۹۰ ضربان در دقیقه) به‌وضوح زیر خط نرمال قرار گرفته‌اند و در دم
بالا نیز برخی مقادیر بالاتر از ۲۰۰ ضربان در دقیقه کمی بالاتر از خط واقع شده‌اند. این انحراف‌ها
نشان‌دهندهٔ کجی ملایم ناشی از مشاهدات بسیار کم یا بسیار بالا هستند که با تغییرات طبیعی پاسخ
بیماران در تست ورزش سازگار است.

\begin{figure}[H]
	\centering
	\includegraphics[width=0.7\textwidth]{qq_thalach.png}
	\caption{نمودار \lr{Q--Q} برای ویژگی \lr{thalach}.}
	\label{fig:qq-thalach}
\end{figure}

به‌طور کلی، نمودارهای \lr{Q--Q} نشان می‌دهند که هیچ‌یک از ویژگی‌های بررسی‌شده کاملاً از
توزیع نرمال پیروی نمی‌کنند. هر ویژگی در دم‌ها یا کلاً رفتاری با کجی و انحراف از نرمال دارد
و ویژگی \lr{chol} بیشترین انحراف را به دلیل وجود مقادیر بسیار بالای کلسترول نشان می‌دهد.
این مسئله برای طبقه‌بند \lr{KNN} مشکل‌زا نیست، زیرا این الگوریتم فرضی دربارهٔ نرمال بودن
توزیع ویژگی‌ها ندارد.

\subsection{ماتریس همبستگی}

برای بررسی روابط خطی میان ویژگی‌های عددی و نیز ارتباط آن‌ها با متغیر هدف، ماتریس همبستگی پیرسون
محاسبه و به صورت یک نقشهٔ حرارتی نمایش داده شد (شکل \ref{fig:corr-matrix}). این ماتریس الگوهای
بالینی معناداری را نشان می‌دهد و مشخص می‌کند کدام ویژگی‌ها برای تمایز بین بیماران دارای و فاقد
بیماری قلبی مهم‌تر هستند.

\begin{figure}[H]
	\centering
	\includegraphics[width=0.75\textwidth]{corr_matrix.png}
	\caption{ماتریس همبستگی برای تعدادی از ویژگی‌های عددی و متغیر هدف.}
	\label{fig:corr-matrix}
\end{figure}

قوی‌ترین همبستگی‌ها با متغیر هدف برای \lr{thalach} و \lr{oldpeak} مشاهده می‌شوند. ویژگی
\lr{thalach} (حداکثر ضربان قلب) همبستگی مثبت متوسطی با هدف دارد ($r = 0.42$) که نشان می‌دهد
افرادی که در طول تست ورزش به ضربان قلب بالاتری می‌رسند، کمتر در معرض ابتلا به بیماری قلبی قرار
دارند. در مقابل، \lr{oldpeak} همبستگی منفی نسبتاً قوی با متغیر هدف دارد ($r = -0.44$) که با
دانش بالینی موجود هم‌خوان است؛ یعنی افت بیشتر قطعهٔ \lr{ST} معمولاً با ایسکمی میوکارد و احتمال
بیشتر بیماری قلبی همراه است. سن نیز همبستگی منفی ضعیف‌تری با هدف دارد ($r = -0.22$) که بیانگر
افزایش شیوع بیماری قلبی با افزایش سن است، هرچند این رابطه کاملاً خطی نیست.

سایر ویژگی‌ها مانند \lr{trestbps} و \lr{chol} تنها همبستگی‌های ضعیفی با متغیر هدف نشان می‌دهند؛
این موضوع حاکی از آن است که در چارچوب این مجموعه‌داده، این متغیرها به‌تنهایی قدرت پیش‌بینی
محدودی دارند، اما ممکن است در ترکیب با سایر ویژگی‌ها در یک فضای چندمتغیره نقش مؤثری ایفا کنند.

برخی روابط میان ویژگی‌ها نیز قابل توجه‌اند. سن و \lr{thalach} همبستگی منفی متوسطی دارند
($r = -0.38$) که نشان می‌دهد افراد مسن‌تر معمولاً در تست ورزش به حداکثر ضربان قلب پایین‌تری
می‌رسند. همچنین، \lr{oldpeak} و \lr{thalach} نیز همبستگی منفی ($r = -0.35$) نشان می‌دهند که
کاهش تحمل ورزش و عملکرد قلبی ضعیف‌تر در بیماران با پاسخ ایسکمیک را منعکس می‌کند.

در مجموع، ماتریس همبستگی نشان می‌دهد که در این مجموعه‌داده، بیماری قلبی بیش از همه با ویژگی‌های
مرتبط با عملکرد در تست ورزش مانند \lr{thalach} و \lr{oldpeak} مرتبط است، در حالی که اندازه‌گیری‌های
استراحتی مانند \lr{trestbps} و \lr{chol} وابستگی خطی ضعیف‌تری با متغیر هدف دارند. این یافته‌ها
همچنین توضیح می‌دهند که چرا مدل‌های مبتنی بر فاصله مانند \lr{KNN}، که تعامل هم‌زمان چند ویژگی را
در نظر می‌گیرند، می‌توانند علیرغم ضعف همبستگی خطی تک‌ویژگی، عملکرد مناسبی داشته باشند.

\subsection{نمودارهای جفتی (\lr{Pair Plots})}

برای بررسی روابط چندمتغیره بین ویژگی‌های عددی کلیدی، یک \lr{pair plot} برای ویژگی‌های \lr{age}،
\lr{chol}، \lr{thalach} و \lr{oldpeak} رسم شد که در آن نقاط بر اساس متغیر هدف دودویی رنگ‌آمیزی شده‌اند.
این نمودار نمایی هم‌زمان از توزیع حاشیه‌ای هر ویژگی و روابط دو به دوی آن‌ها ارائه می‌دهد
(شکل \ref{fig:pairplot-selected}).

\begin{figure}[H]
	\centering
	\includegraphics[width=0.85\textwidth]{pairplot_selected.png}
	\caption{نمودار جفتی برای چند ویژگی عددی منتخب، رنگ‌آمیزی‌شده بر اساس کلاس هدف.}
	\label{fig:pairplot-selected}
\end{figure}

هیستوگرام‌های روی قطر نشان می‌دهند که \lr{age} و \lr{thalach} توزیع‌هایی نسبتاً متقارن دارند، در حالی
که \lr{chol} و \lr{oldpeak} به‌طور قابل توجهی به سمت راست کج شده‌اند. از نظر توزیع کلاس‌ها، بیماران
دارای بیماری قلبی (هدف = ۱) اندکی بیشتر در محدودهٔ سن‌های بالاتر دیده می‌شوند، هرچند این اثر نسبتاً
ضعیف است. در مقابل، تفاوت‌های واضح‌تری در توزیع \lr{thalach} و \lr{oldpeak} مشاهده می‌شود: افرادی که
بیماری قلبی ندارند معمولاً به حداکثر ضربان قلب بالاتر می‌رسند، در حالی که بیماران قلبی بیش‌تر
\lr{oldpeak}های بزرگ‌تر و \lr{thalach}های کم‌تر دارند.

نمودارهای پراکندگی دو به دو الگوهای مهمی را آشکار می‌کنند. رابطهٔ بین \lr{age} و \lr{thalach}
شیب منفی مشخصی دارد که کاهش ظرفیت ورزش با افزایش سن را نشان می‌دهد. همچنین، ترکیب \lr{thalach} و
\lr{oldpeak} بهترین جداسازی بصری بین دو کلاس را فراهم می‌کند: افراد بدون بیماری قلبی عمدتاً در
ناحیهٔ \lr{thalach} بالا و \lr{oldpeak} نزدیک صفر خوشه‌بندی می‌شوند، در حالی که بیماران قلبی
در مقادیر پایین‌تر \lr{thalach} و مقادیر بالاتر \lr{oldpeak} متمرکز هستند. در مقابل، ویژگی \lr{chol}
در تمام نمایش‌های دو به دو هم‌پوشانی قابل‌توجهی بین دو کلاس نشان می‌دهد و حاکی از آن است که این
ویژگی به تنهایی توان تفکیک‌کنندگی ضعیف‌تری دارد.

در مجموع، نمودارهای جفتی تأیید می‌کنند که بیماری قلبی در این مجموعه‌داده بیش از همه با شاخص‌های
عملکردی مرتبط با تست ورزش (\lr{thalach} و \lr{oldpeak}) پیوند دارد، در حالی که ویژگی‌هایی مانند
\lr{chol} و \lr{age} جداسازی ضعیف‌تر یا کم‌ثبات‌تری ارائه می‌دهند. این الگوهای چندمتغیره استفاده
از طبقه‌بند مبتنی بر فاصله‌ای مانند \lr{KNN} را توجیه می‌کنند، زیرا چنین مدلی می‌تواند از این
روابط حتی در صورت غیرخطی بودن یا نبود جداسازی واضح در امتداد یک ویژگی منفرد بهره‌برداری کند.

\subsection{نمودارهای شمارشی ویژگی‌های دسته‌ای}

برای درک بهتر توزیع ویژگی‌های دسته‌ای در مجموعه‌داده، نمودارهای شمارشی برای \lr{cp} (نوع درد
قفسهٔ سینه)، \lr{exang} (آنژین ناشی از ورزش)، \lr{sex} و متغیر هدف \lr{target} رسم شد. این
نمودارها فراوانی هر دسته را نشان می‌دهند و به شناسایی عدم توازن‌های احتمالی در داده کمک می‌کنند.

\paragraph{نوع درد قفسهٔ سینه (\lr{cp}).}
شکل \ref{fig:count-cp} نشان می‌دهد که توزیع نوع درد قفسهٔ سینه یکنواخت نیست. دستهٔ
\lr{cp = 0} (آنژین تیپیک) با نزدیک به ۵۰۰ مشاهده، رایج‌ترین حالت است، در حالی که
\lr{cp = 3} (بدون علامت) با کمتر از ۱۰۰ نمونه کم‌ترین فراوانی را دارد. دو دستهٔ دیگر از نظر
فراوانی بین این دو قرار می‌گیرند. این الگو با مشاهدات بالینی واقعی سازگار است، زیرا آنژین
تیپیک اغلب شکایت اصلی در بیماران قلبی محسوب می‌شود.

\begin{figure}[H]
	\centering
	\includegraphics[width=0.7\textwidth]{count_cp.png}
	\caption{نمودار شمارشی ویژگی \lr{cp} (نوع درد قفسهٔ سینه).}
	\label{fig:count-cp}
\end{figure}

\paragraph{آنژین ناشی از ورزش (\lr{exang}).}
همان‌طور که در شکل \ref{fig:count-exang} دیده می‌شود، اکثریت بیماران (\lr{exang = 0}) در طول
تست ورزش دچار آنژین نمی‌شوند. زیرمجموعهٔ کوچک‌تری، تقریباً حدود یک‌سوم داده‌ها، شامل بیمارانی است
که در طول تست دچار آنژین می‌شوند (\lr{exang = 1}). این عدم توازن تا حدی طبیعی است، زیرا همهٔ افراد
با وجود بیماری زمینه‌ای الزاماً در تست ورزش دچار درد قفسهٔ سینه نمی‌شوند.

\begin{figure}[H]
	\centering
	\includegraphics[width=0.7\textwidth]{count_exang.png}
	\caption{نمودار شمارشی ویژگی \lr{exang}.}
	\label{fig:count-exang}
\end{figure}

\paragraph{جنسیت (\lr{sex}).}
شکل \ref{fig:count-sex} نشان می‌دهد که در این مجموعه‌داده تعداد بیماران مرد (\lr{sex = 1})
به‌طور قابل‌توجهی بیشتر از بیماران زن (\lr{sex = 0}) است و نسبت تقریبی ۲ به ۱ را نشان می‌دهد.
این الگو با شیوع بالاتر بیماری قلبی در مردان در بسیاری از جمعیت‌های بالینی هم‌خوانی دارد.

\begin{figure}[H]
	\centering
	\includegraphics[width=0.7\textwidth]{count_sex.png}
	\caption{نمودار شمارشی ویژگی \lr{sex}.}
	\label{fig:count-sex}
\end{figure}

\paragraph{متغیر هدف (\lr{target}).}
همان‌طور که در شکل \ref{fig:count-target} مشاهده می‌شود، توزیع کلاس در متغیر هدف تقریباً
متعادل است؛ به این معنا که تعداد نمونه‌های کلاس \lr{target = 0} (بدون بیماری قلبی) و
\lr{target = 1} (وجود بیماری قلبی) تقریباً برابر است. این تعادل برای آموزش مدل‌های
طبقه‌بندی مفید است، زیرا ریسک سوگیری مدل به سمت یکی از کلاس‌ها را کاهش می‌دهد و نیاز به
روش‌های متعادل‌سازی مانند oversampling یا وزن‌دهی به کلاس‌ها را برطرف می‌کند.

\begin{figure}[H]
	\centering
	\includegraphics[width=0.7\textwidth]{count_target.png}
	\caption{نمودار شمارشی متغیر هدف \lr{target}.}
	\label{fig:count-target}
\end{figure}

در مجموع، نمودارهای شمارشی نشان می‌دهند که توزیع ویژگی‌های دسته‌ای در این مجموعه‌داده با الگوهای
کلینیکی واقعی سازگار است. اگرچه برخی دسته‌ها، مانند نوع درد قفسهٔ سینه و جنسیت، عدم توازن قابل‌توجهی
دارند، اما توازن مناسب در متغیر هدف باعث می‌شود مدل‌هایی مانند \lr{KNN} بدون نیاز به تنظیمات
اضافی برای عدم توازن کلاس‌ها، به‌طور مؤثر آموزش داده شوند.